\documentclass[11pt,a4paper]{article}

\usepackage[margin=2.5cm]{geometry}
\usepackage[T1]{fontenc}
\usepackage{lmodern}
\usepackage{amsmath}
\usepackage{booktabs}
\usepackage{tabularx}
\usepackage{enumitem}
\usepackage{xcolor}
\usepackage{listings}
\usepackage{hyperref}

\hypersetup{colorlinks=true, linkcolor=blue!50!black, urlcolor=blue!50!black}
\lstset{basicstyle=\ttfamily\small, columns=fullflexible, breaklines=true, frame=single,
  framerule=0.3pt, rulecolor=\color{black!30}, backgroundcolor=\color{black!3}, showstringspaces=false}
\newcommand{\code}[1]{\texttt{#1}}

\title{CPU--GPU Co-Scheduling Module for CloudSim 7G\\[4pt]
\large Phase 1c Report: File Formats, Performance Models and Workload}
\author{Chirag \and Hamiz}
\date{22 November 2026}

\begin{document}
\maketitle

\begin{abstract}
Phase~1c builds the inputs of a simulation run. It fixes the three file formats (job CSV, profile
JSON, experiment YAML) shared with the Python tools in \code{tools/}; loads Synergy's 69 measured
throughput profiles, converted by \code{tools/convert\_synergy\_profiles.py}; implements the throughput,
locality and rate models; and implements the workload reader and model labellers. Two results stand
out. First, the class-split labeller reproduces Synergy's model label for all 1,000 measured jobs of
the reference run, so validation can regenerate labels instead of copying them. Second, Synergy's
Fig.~12 profiles keep one throughput table but state speedups relative to each cluster's own fixed
share; normalising to the server's fixed share reproduces them. All 31 tests of the co-scheduling
module pass.
\end{abstract}

\tableofcontents

%======================================================================
\section{Inputs from the tools}
%======================================================================

The tools in \code{tools/} (profile converter, Synergy job extractor, Philly preparation, run
comparator) were added in commit \code{06a3cc98}. They are the only code that reads raw traces or
Synergy's files; the simulator reads only the formats below. The converter was checked against
Synergy's source before its output was used:

\begin{itemize}
  \item Rows of Synergy's \code{tput} matrix are cores per GPU $\{1,2,3,4,5,6,9,12,24\}$ and columns
        are GB of memory per GPU $\{20.83, 62.5, 125, 187.5, 250\}$ (\code{jobs/job.py}, lines 96--97).
  \item Synergy's simulator loads only \code{<model>\_1} files (\code{jobs/model.py}, line 171).
  \item ResNet-18 gives 1.0 at 3 cores and 62.5\,GB per GPU (the fixed share on a 24-CPU, 8-GPU,
        500\,GB server) and 2.4 at 9 cores, as stated in the architecture document.
\end{itemize}

The converter's own tests pass (26 passed; one is skipped because it needs a Synergy \code{out.log}).
The converted profiles are stored in \code{presets/profiles/synergy-v100/} with their source and
licence (MIT).

%======================================================================
\section{File formats}
%======================================================================

The formats are documented in \code{docs/formats.md}.

\paragraph{Job CSV.} Required: \code{job\_id}, \code{arrival\_s}, \code{gpus}, and work as either
\code{iterations} (Synergy-style, optionally with \code{iter\_time\_s}) or \code{duration\_s}
(trace-style: run time at the fixed share). Optional: \code{model}, \code{team}, \code{priority}.
Other columns are ignored, so a reference file with \code{jct\_s} can be read as a workload. Every
error names its line, and duplicate job ids are rejected.

\paragraph{Profile JSON.} One file per (model, GPU type, GPU count). The Java loader reads the
fields the converter writes, ignores the rest, and validates the axes (strictly increasing) and the
table shape.

\paragraph{Experiment YAML.} Cluster, organisation, workload and labeller, model settings, the
candidate schedulers (each a trigger, ordering, allocation, placement and lifecycle plug-in with
parameters), seeds and output folder. Unknown keys are errors; relative paths are relative to the
YAML file.

%======================================================================
\section{Performance models (L2)}
%======================================================================

\subsection{Profile repository}
Profiles are looked up by (model, GPU type, GPU count). With \code{scaling: linear} (what Synergy
does) a job on any number of GPUs uses the 1-GPU profile; with \code{per\_gpu\_count} it uses the
profile measured for its exact GPU count. A missing profile stops the run with a list of what is
available; there is never a fallback to another GPU type. \code{requireAll} checks every needed
combination before a run starts.

\paragraph{Unmeasured points.} The multi-GPU files (not used by Synergy) contain all-zero rows at
high core counts: 12 and 24 cores per GPU for 2 GPUs, 9 to 24 for 4 GPUs. These mean ``not measured''
(for 4 GPUs, 9 cores per GPU would need 36 of a server's 24 cores). The loader accepts only trailing
all-zero rows and treats the last non-zero row as the largest measured point.

\subsection{Throughput model}
Speedup is always relative to the job's speed at the fixed share of the server it runs on:
\[
  \text{speedup}(c, m) = \frac{T(c, m)}{T(c_\text{fair}, m_\text{fair})}, \qquad
  c_\text{fair} = \frac{\text{server cores}}{\text{server GPUs}}, \quad
  m_\text{fair} = \frac{\text{server memory}}{\text{server GPUs}},
\]
where $c$ and $m$ are the job's cores and memory per GPU and $T$ is the profile table. Between
measured CPU points $T$ is linearly interpolated; the memory column is the largest measured value at
or below $m$. Below the first point the first value is used. Beyond the last measured point $T$ stays
at the last value and a warning is recorded once per model. Nothing is extrapolated.

\paragraph{Finding: the Fig.~12 profiles.} Synergy's \code{models-4}, \code{models-5} and
\code{models-6} belong to servers with 32, 40 and 48 CPUs. Their tables are identical to the default
ones (1.0 at 3 cores per GPU), but their \code{ideal\_speedup} is lower: for ResNet-18, 1.85, 1.568
and 1.37 instead of 2.4. These equal $T(9)/T(4)$, $T(9)/T(5)$ and $T(9)/T(6)$, the speedup over each
cluster's own fixed share of 4, 5 and 6 cores per GPU. This holds for 46 of the 48 single-GPU
profiles. Normalising to the server's fixed share, as above, therefore reproduces Synergy's values
from the table alone; a model that always divided by the 3-core value would be wrong on these
clusters.

The other two profiles, \code{models-6/alexnet} and \code{models-6/shufflenet}, have a changed
\code{ideal\_speedup} and \code{iter\_time} but an unchanged table, so table and speedup disagree in
Synergy's own data. This is recorded in the profiles' README.

\paragraph{Two modes.} \code{table} uses the formula above. \code{synergy\_two\_point} does exactly what
Synergy's simulator does: a job holding at least its ideal cores and memory per GPU runs at
\code{ideal\_speedup}, anything else at 1.0. Synergy's Greedy and Tune only ever allocate these two
points, so validation presets use this mode and are exact even for the two inconsistent profiles.

\paragraph{Ideal demand and GPU busy fraction.} The ideal demand is the profile's ideal cores and
memory per GPU times the job's GPUs, capped at one server. The GPU busy fraction is the speedup at the
current allocation divided by the speedup at the ideal demand, at most 1; for ResNet-18 at the fixed
share this is $1/2.4 \approx 42\%$. It drives GPU power in the GPU layer.

\subsection{Locality, interference and rate}
\code{LocalityModel} returns 1 for a job on one server and the penalty $p$ when it is spread
(validation: $p = 1$, as in all of Synergy's active profiles). \code{InterferenceModel} is the named
seam for a future measured slowdown; in this version it is 1. \code{RateCalculator} is the single
place a job's speed is computed, following Synergy's \code{jobs/job.py}:
\[
  \text{rate} = \frac{\text{speedup} \times \text{locality} \times \text{interference}}{\text{iteration time at the fixed share}}
  \quad \text{iterations per second.}
\]
The iteration time is the job's own when the job CSV gives one (the reference run does, per job) and
the profile's otherwise.

%======================================================================
\section{Workload (L3)}
%======================================================================

\code{JobCsvReader} reads the job CSV; \code{TraceWorkload} orders jobs by arrival (ties by id) and
labels each one in that order. Two labellers are implemented.

\paragraph{From trace.} Uses the CSV's \code{model} column; a job without one is an error.

\paragraph{Class split (Synergy's rule).} Synergy's simulator uses assignment mode \code{OVERALL}
(\code{jobs/workload.py}, line 62). Each class (image, language, speech) has a priority
\[
  \text{priority} = \frac{\text{target percent}}{100} - \frac{\text{jobs of the class so far}}{\text{all jobs so far}},
\]
the job goes to the class with the highest priority (ties to the earlier class, as Python's stable
sort does), and within the class models are taken round-robin: image \{alexnet, res18, res50,
mobilenet, shufflenet\}, language \{gnmt, transformer, lstm\}, speech \{m5, deepspeech\}.

\paragraph{Finding: labels are reproducible.} Labelling jobs 0 to 4999 in order with split
(20, 70, 10) gives exactly Synergy's model for all 1,000 measured jobs (4000--4999) of the reference
run, and the expected class totals (1000, 3500, 500). The architecture document assumed that labels
depend on which jobs are running and therefore cannot be regenerated; that is true only of Synergy's
unused \code{RUNNABLE} mode. In \code{OVERALL} mode labels depend only on arrival order.

%======================================================================
\section{Tests}
%======================================================================

\begin{center}
\begin{tabularx}{\linewidth}{@{}lX@{}}
\toprule
Test & Checks \\
\midrule
fidelity & Every single-GPU profile Synergy uses: 1.0 at the fixed share and \code{ideal\_speedup} at its
  ideal demand (table mode within 0.02; two-point mode exactly). \\
ResNet-18 & 2.4 at 9 cores per GPU; ideal demand of a 2-GPU job is 18 cores, 125\,GB. \\
Fig.~12 & On a 32-CPU server with \code{models-4}: 1.0 at 4 cores per GPU, 1.85 at the ideal. \\
interpolation & 7.5 cores per GPU is half way between the 6- and 9-core values; memory below 62.5\,GB
  uses the 20.83\,GB column. \\
beyond measured & Held at the last value; one warning per model. \\
missing profile & Unknown GPU type, unknown model and a missing GPU count all stop the run. \\
scaling & \code{linear} picks the 1-GPU profile for 4 GPUs; \code{per\_gpu\_count} the 4-GPU one. \\
locality, rate & Penalty only when spread; rate $= 2.4/0.64$ at the ideal, $\times 0.8$ when spread,
  $1/0.64$ at the fixed share with GPUs 42\% busy. \\
job CSV & Both work styles, extra columns ignored, errors with line numbers, duplicates rejected. \\
labellers & From-trace needs a model; class split reproduces Synergy's 1,000 labels and starts with the
  first class and round-robins within classes. \\
YAML & The documented example parses; unknown keys and incomplete candidates are rejected. \\
\bottomrule
\end{tabularx}
\end{center}

Co-scheduling module: 31 tests, 0 failures (\code{mvn install}).

%======================================================================
\section{Next step}
%======================================================================
Phase~2 connects these inputs to the GPU layer: the policy interfaces, FIFO with GPU-proportional
allocation and consolidated placement, the broker that builds a round plan and turns it into GPU-layer
commands, per-job metrics, the run manifest and the experiment runner, and the first end-to-end run.

\end{document}
